\section{Mergeable Sketches, Neural Operators, and C-TreePO}
\label{sec:mergeable-sketches}

A query system processes events split across many machines and asks a
global question (how many distinct users? what is the median latency?)
without holding the full data anywhere. The data-systems literature
handles this with \emph{mergeable summaries}
\citep{Agarwal2013MergeableTODS}: each shard (the classical analogue
of a C-Tree leaf block) is replaced by a compact state, the states
combine pairwise up a tree, and the query is answered once, at the
root. Closely related strands include database algebraic
aggregates \citep{GrayEtAl1997DataCube}, distributed-streaming work on
the symmetric unordered case \citep{FeldmanEtAl2008MUD}, and free-monoid
homomorphisms in functional programming
\citep{Gibbons1996ThirdHomomorphism}. C-TreePO takes this lineage as its
state-level algebraic target: learn a bounded state, merge it locally,
and audit whether the task answer is preserved. The summarizer \(g\)
plays the role of the encode function at leaves and the merge operator
at internal nodes. The classical query at the root becomes a learned
query \(f\) on the stored state. An oracle \(\fstar\) names the task
value the researcher cares about (an expert label, a calibrated reward,
a preference judgment), and the learned query is calibrated against it.
Three local laws C1, C2, C3 check that the learned summarizer preserves
enough information at each step for the root answer to track the
oracle. The same residuals are checked twice: training data drive them
as a soft loss that pushes \(g\) toward local-law compatibility, and a
runtime audit then estimates them on the realized states a tree
produces, converting them into a certificate that replaces the
algebraic correctness proof a classical sketch would supply.
Section~\ref{sec:framework} develops the formal versions.

Child answers alone can leave the parent answer ambiguous. State-level
merging is what makes hierarchical aggregation possible, and the local
laws are how a learned state earns the same privilege.

\paragraph{Notation used in this section.}
A handful of symbols recur throughout. The data living at a tree node
\(v\) is \(X_v\in\mathcal{X}\); at a leaf this is just the raw block
\(b_v\). Leaf blocks live in a leaf-object class
\(\mathcal{B}\subseteq\mathcal{X}\), and produced states live in a
state class \(\mathcal{Z}\subseteq\mathcal{X}\), after rendering if
necessary. The state stored at \(v\) is \(s_v\). A query \(f\) reads an
answer off any queryable object in \(\mathcal{X}\), and two objects are
\emph{query-equivalent}, written \(u\sim_f v\), when the query gives
them the same answer. The summarizer \(g\) produces and combines
states; \(\oplus\) is the interface \(g\) uses to combine them, and
\(\uplus\) is multiset addition on the underlying data.
Table~\ref{tab:v7-section2-notation} consolidates these.

\begin{table}[t]
\centering
\small
\begin{tabular}{@{}p{0.24\linewidth}p{0.68\linewidth}@{}}
\toprule
Symbol & Meaning \\
\midrule
\(D\), \(D_1\uplus D_2\) &
Classical dataset or stream multiset; \(\uplus\) is multiset addition,
so repeated records remain repeated. \\
\(X_v\in\mathcal X\) &
The queryable object represented by tree node \(v\). At a leaf,
\(X_v=b_v\); internal nodes represent the fixed composition of their
children. \\
\(\mathcal B\subseteq\mathcal X\), \(\mathcal Z\subseteq\mathcal X\) &
Leaf-block objects and produced summary states, viewed as queryable
objects. In typed implementations a renderer
\(\rho:\mathcal Z\to\mathcal X\) supplies the inclusion for states. \\
\(s_v\), \(z\), \(\mathcal Z\) &
The realized state at node \(v\), a generic state, and the state space. \\
\(g\) &
The summarizer, used to summarize leaves, merge child states, and
re-summarize produced states for C2. \\
\(f\), \(\fstar\) &
The query applied to states or represented objects, and the oracle
value the query is calibrated against. \\
\(u\sim_f v\), \(u\sim_{f,\varepsilon}v\) &
Exact and \(\varepsilon\)-approximate query equivalence. The shorthand
\(u\approx v\) in figures denotes the same relation. \\
\(\oplus\), \(\oplus_g\) &
Classical state merge, or the fixed ordered interface passed to \(g\).
In the simplest case \(\oplus\) is just concatenation of states (the
\emph{identity merge}, where combining two states is left entirely to
\(g\)); the induced learned merge is then
\(z\oplus_g z'=g(z\oplus z')\). \\
\bottomrule
\end{tabular}
\caption{Notation reminder for \S2. The prose above introduces each
symbol in plain language; this table consolidates them.}
\label{tab:v7-section2-notation}
\end{table}

\paragraph{Type convention.}
Throughout this section, \(\mathcal X\) is the queryable object space.
We view leaf blocks and produced states as queryable by writing
\[
  b_v\in\mathcal B\subseteq\mathcal X,\qquad
  s_v,z\in\mathcal Z\subseteq\mathcal X,\qquad
  f:\mathcal X\to\Y .
\]
Thus \(f(b_v)\), \(f(s_v)\), and \(f(X_v)\) are all meaningful. If an
implementation keeps states in a separate type, the renderer
\(\rho:\mathcal Z\to\mathcal X\) is applied before \(f\); below we
suppress \(\rho\) and write \(f(z)\) directly.

The two operators \(\uplus\) and \(\oplus\) sit on different layers of
the same hierarchy. \(\uplus\) combines the underlying data: two
shards, or equivalently the leaves under two subtrees, become one
stream multiset. \(\oplus\) combines the objects produced from that
data, whether two sketch states, two rendered summaries, or the
ordered child package handed to \(g\). The two are linked by the
mergeable-summary guarantee: if \(z_i\) summarizes \(D_i\), then
\(z_1\oplus z_2\) summarizes \(D_1\uplus D_2\), so a state-level merge
represents the data-level merge it stands in for. For ordered C-Trees
the analogous object-level composition is \(X_L\oplus X_R\); it
preserves order and is therefore distinct from multiset addition.

\subsection{Mergeable Summaries: Definition and Requirements}
\label{sec:v7-mergeable-pattern}

A \emph{mergeable summary} is a compact data structure equipped with a
binary merge operation that behaves well under hierarchical aggregation.
Following \citet{Agarwal2013MergeableTODS}, a summary scheme for a query
family \(\mathcal{Q}\) (cardinality, frequency, quantile, membership,
or any other aggregate the system needs to compute) has three
components: an \emph{encode} function that turns a shard \(D\) (a
C-Tree leaf block in our language) into a state, a \emph{merge}
operator \(\oplus\) that combines two states, and a \emph{query}
function that reads the answer off a state. The scheme is mergeable
when, for any two shards or leaf blocks \(D_1,D_2\), both
\(\mathrm{encode}(D_1\uplus D_2)\) and
\(\mathrm{encode}(D_1)\oplus\mathrm{encode}(D_2)\) answer the query on
the combined leaf set equally well, with state size bounded by some
\(k\) after merging. Operationally, a MapReduce, Dataflow, or Spark job
can encode each shard on a separate worker. In C-Tree language, those
shards are leaf blocks. The system then merges leaf states pairwise up
the tree and reads the answer off the root, without materializing the
full data.

\paragraph{Types and what the state carries.}
Each component has a fixed type:
\[
  \mathrm{encode}:\Strings\to\mathcal{S},\qquad
  \oplus:\mathcal{S}\times\mathcal{S}\to\mathcal{S},\qquad
  \mathrm{query}:\mathcal{S}\times\mathcal{Q}\to\mathcal{R}.
\]
The encode function lives at leaves; the merge \(\oplus\) lives at
internal nodes; the query is asked once, at the root. Each shard or
leaf block's state has to carry whatever is needed to combine with
another shard's state, or another leaf block's state, and still answer
\(Q\) correctly. Distinct count makes this concrete. Knowing how many
distinct items are in leaf block \(A\) and how many in leaf block \(B\)
does not determine how many are in \(A\uplus B\), because the overlap
is unknown. A distinct-count summary therefore stores enough per-leaf
information (typically a hash-based fingerprint of the items seen) for
the merge to detect overlap. Different schemes carry different per-leaf
information, calibrated to what their query needs.

\paragraph{Three properties of the merge operator.}
Three algebraic properties of \(\oplus\) shape what the scheme buys.

\emph{Associativity} means \((a\oplus b)\oplus c = a\oplus(b\oplus c)\).
The schedule of merges does not matter; balanced trees, unbalanced
trees, and streaming reductions all produce the same state. This is
what makes parallel computation possible.

\emph{Commutativity} means \(a\oplus b = b\oplus a\). The order of
inputs does not matter. Appropriate when the data is genuinely
unordered (a multiset of records); inappropriate when order carries
meaning (text, time series, dialogue).

\emph{Idempotence} means \(a\oplus a = a\). Re-merging a state with
itself does not drift. Holds for set- or semilattice-like unions
(HyperLogLog, Bloom filters); fails for additive counter sketches like
Count-Min, where re-merging would double the counts.

For ordered data such as text, the relevant algebraic structure is a
free-monoid homomorphism\footnote{A \emph{free monoid} over an alphabet
\(\Sigma\) is the set of finite sequences over \(\Sigma\) under
concatenation, with the empty sequence as identity. A
\emph{homomorphism} into another monoid \((M,\odot)\) is a map \(h\)
that respects concatenation: for sequences \(u\) and \(v\),
\(h(u\oplus v)=h(u)\odot h(v)\), and it sends the empty sequence to
the identity of \(\odot\).}
\citep{Gibbons1996ThirdHomomorphism}: the summarizer respects
concatenation but not reordering, \(h(u\oplus v)=h(u)\odot h(v)\),
with an associative \(\odot\) that need not be commutative.

\paragraph{Validity, size profile, mergeability.}
What counts as a valid state depends on two ingredients: the input being
summarized and the target accuracy. The input \(D\) is a dataset or
stream multiset, in which a record appearing twice contributes two
copies, and the target accuracy is a number \(\varepsilon\) that bounds
the allowable error of the final query. We write \(S(D,\varepsilon)\)
for any state produced for \(D\) at accuracy \(\varepsilon\). The
notation is one-to-many because different stream orders or random
choices may give different valid states for the same input.

\begin{definition}[Valid summary state]
\label{def:v7-valid-state}
A state \(z=S(D,\varepsilon)\) is \emph{valid} when:
\begin{enumerate}[leftmargin=1.5em,nosep]
\item the type of \(z\) matches the scheme's state interface
(well-formedness);
\item \(|z|\) is within the advertised size bound \(k\);
\item the query answered from \(z\) has error at most \(\varepsilon\).
\end{enumerate}
\end{definition}

A valid state is the kind of compact summary a distributed system is
allowed to keep: it has the right format, fits the advertised memory
budget, and answers the target query accurately enough. Each clause
does distinct work, and each connects to a separate piece of the
scheme. The size bound is what makes the state \emph{compact};
tracking how the bound grows with input size yields the \emph{size
profile}
\[
  k(n,\varepsilon)
  :=
  \max\{\,|S(D,\varepsilon)| : |D|=n,\ S(D,\varepsilon)
    \text{ is valid}\,\},
\]
abbreviated \(k(\varepsilon)\) when independent of \(n\). Uniform
bounds in \(n\) are the useful regime, since a system can reserve a
fixed memory footprint per shard or leaf block regardless of how much
data that block contains. The query-error bound is what makes the state
\emph{useful}
for the target query, and the defining law of a mergeable scheme is
that \(\oplus\) preserves it: whenever \(z_1\) is a valid
\(S(D_1,\varepsilon)\) and \(z_2\) is a valid \(S(D_2,\varepsilon)\),
\[
  z_1 \oplus z_2 \quad\text{is a valid}\quad
  S(D_1\uplus D_2,\varepsilon),
\]
where \(D_1\uplus D_2\) is multiset addition (repeated records remain
repeated). The well-formedness clause is what gives \(\oplus\) a typed
object to act on. A tree of merges built from validly-summarized leaves
therefore produces a valid root state, and the query is asked once, at
that root.

\paragraph{Mergeability for general object spaces.}
Mergeable summaries lift from record multisets to whatever object
space a tree actually carries. C-Trees may represent their inputs as
one-hot tokens, supplied embeddings, JSON state, or rendered text;
the mergeability requirement does not depend on which representation
is chosen, only that \(g\) and \(f\) operate consistently on it. Let
\(\mathcal{X}\) be the chosen object space and \(X\in\mathcal{X}\)
the C-Tree representation of \(D\). Fix \(\varepsilon\) for the rest
of the discussion, and let \(\mathcal{Z}\) be the state space and
\(\mathrm{Valid}(\cdot,\cdot)\) the validity relation at that
accuracy. A mergeable summary at this level is a state
\(z\in\mathcal{Z}\) with \(\mathrm{Valid}(X,z)\): well formed, within
the size bound, and preserving the root task value to accuracy
\(\varepsilon\). Mergeability lifts in the natural way: if two states
are valid summaries of two objects in \(\mathcal{X}\), their merge is
a valid summary of the parent object the tree composes from them. Let
\(X_{12}\) denote that parent object, obtained by the tree's fixed
composition rule. Formally,
\[
  \mathrm{Valid}(X_1,z_1)
  \wedge
  \mathrm{Valid}(X_2,z_2)
  \Longrightarrow
  \mathrm{Valid}(X_{12},\,z_1\oplus z_2).
\]
The display is the same statement as ``\(z_1\oplus z_2\) is a valid
\(S(D_1\uplus D_2,\varepsilon)\)'' above, with the dataset multiset
\(D\) replaced by the represented parent object \(X_{12}\) and the
one-to-many summary \(S(\cdot,\varepsilon)\) replaced by the relation
\(\mathrm{Valid}(\cdot,\cdot)\). Table~\ref{tab:v7-agarwal-translation}
records the column-by-column correspondence.

The operational mapping to C-Trees is direct. A classical shard
\(D_i\) is represented by a C-Tree leaf block \(b_i\) or by the
queryable object \(X_i\) attached to that leaf. Classical
\(\mathrm{encode}(D_i)\) becomes the leaf state \(s_i=g(b_i)\).
Classical state merge \(z_L\oplus z_R\) becomes the internal C-Tree
call \(g(s_L\oplus s_R)\). The classical final query becomes the
query \(f(s_{\mathrm{root}})\), applied only after all leaf states
have been merged to the root.

\begin{table}[t]
\centering
\small
\setlength{\tabcolsep}{4pt}
\begin{tabular}{@{}p{0.28\linewidth}p{0.34\linewidth}p{0.28\linewidth}@{}}
\toprule
Mergeable-summary notation & State-level notation & C-Tree role \\
\midrule
\(S()\) & summary method / state interface & choice of state space, validity, merge, and query \\
\(S(D,\varepsilon)\) & a state \(z\) with
\(\mathrm{Valid}(X,z)\) & a valid summary state for the leaf block or
subtree object \(X\) corresponding to \(D\) \\
\(z_1\oplus z_2\) at the merge step & \(z_1\oplus z_2\) & internal \(g\)-call on child states \\
\(D_1\uplus D_2\) & represented parent object \(X_{12}\) & data represented by the parent subtree, obtained by joining the corresponding leaves \\
\(k(n,\varepsilon)\) & maximum size of valid states for \(|D|=n\) & state-size profile preserved by the tree \\
final query on \(S(D,\varepsilon)\) &
\(\mathrm{query}(z,\cdot)\) & root answer \(f(s_{\mathrm{root}})\) after specializing the query family \\
\bottomrule
\end{tabular}
\caption{Mapping between
\citet{Agarwal2013MergeableTODS}'s summary notation and the state-level
notation used in the C-Tree correspondence, after fixing
\(\varepsilon\).}
\label{tab:v7-agarwal-translation}
\end{table}

\paragraph{Variants.}
\emph{One-way mergeability} is the streaming special case: one side of
the merge is a fresh singleton or raw suffix, so merges are between an
existing state and a new record. Equal-size or restricted-size variants
constrain which pairs of states may be merged.
% TODO(citations): cite specific one-way / restricted-size mergeability
% papers beyond \citet{Agarwal2013MergeableTODS}.

\emph{Randomized validity.} For randomized summaries, validity is an
event. A single random seed \(\omega\in\Omega\) may encode all random
choices used by a merge-tree evaluation. The Agarwal-style success
statement says that, for any merge tree whose leaves represent \(D\),
the root state lies in \(S(D,\varepsilon)\) with probability at least
\(p\). C-Trees inherit this probability shape; see
\S\ref{sec:v7-ctreepo-correspondence} below.
% TODO(citations): cite the specific randomized-mergeable-summary
% treatments (KLL, sampling sketches) used here.

\emph{Range spaces.} For range spaces the subset notation is literal. A
finite range space is \((D,\mathcal{R})\), and an
\(\varepsilon\)-approximation is a subset \(A\subseteq D\) satisfying
\[
  \sup_{R\in\mathcal{R}}
  \left|
    \frac{|A\cap R|}{|A|}
    -
    \frac{|D\cap R|}{|D|}
  \right|
  \le \varepsilon.
\]
% TODO(citations): cite the standard \(\varepsilon\)-approximation /
% VC-dimension references (e.g.\ Vapnik--Chervonenkis, Chazelle) for the
% range-space construction.
The selected subset \(A\) reproduces every range's share of the full
data \(D\) up to error \(\varepsilon\). In the C-Tree translation,
\(X_{12}\) is the represented parent object
tracked by the tree, while \(A\) is the selected subset state.
Keeping those roles separate prevents the mistaken inference that
scalar child answers should themselves merge.

\subsection{From Classical Sketches to Learned \(f\) and \(g\)}
\label{sec:v7-from-classical-to-learned}

The mergeable-summary picture above is classical: a hand-designed state
class \(\mathcal{Z}\), a hand-designed merge, and a closed-form
\(\mathrm{query}\) for one fixed query family. C-TreePO targets
settings where this engineering is not feasible. The query a researcher
actually cares about is a task value (an expert label, a calibrated
reward, a preference judgment), and no hand-designed sketch exists for
it. Even when one exists, it is specific to a single query: a different
task means a different sketch, a different state class, a different
merge.

C-TreePO learns the summarizer and the query for the task at hand.
The summarizer \(g\) plays the role of \(\mathrm{encode}\) at leaves
and induces the learned merge \(\oplus_g\) at internal nodes, as a
neural network or a prompted language model operating on tokens,
embeddings, or rendered text. The query \(f\) plays the role of
\(\mathrm{query}\) at the root, as a learned, calibrated estimate of an
oracle \(\fstar:\mathcal{X}\to\Y\) (the task value the researcher
would assign if they could afford to label every document). One
assumption on the oracle matters from the start: it is taken to be
\emph{locally context-compatible}, meaning that substituting one
passage for another with the same task value, in the same
surrounding text, leaves the task value of the whole unchanged.
This is Assumption~\ref{ass:context}, stated formally in
\S\ref{sec:v7-ctreepo-state-machine}, and it is a real restriction. Tasks whose
oracle is sensitive to discourse (irony, rhetorical tone, anything
where local substitution of equivalent content changes how the
surrounding text is read) are out of scope.

A classical sketch proves correctness once, by the algebra of the
sketch, and never has to re-check at runtime. C-TreePO replaces that
algebraic proof with two complementary pressures on the same
residuals: training data penalize them as part of the loss, and a
runtime audit certifies them on the realized states a deployed tree
produces. Because \(g\) is learned, the local laws below cannot be
derived as identities on the operator; they have to be verified
empirically, and either source of evidence (training data or audit)
operates on the realized states the tree generates. The audit samples
realized residuals on a deployed run (the actual leaf summaries and
merge calls produced for one document) and treats the run as
\emph{realized-valid} (valid on the actual tree calls made in this run)
when the samples support a chosen bound. The
samples come from oracle access where available, or from randomly
drawn root-level supervision otherwise.
Section~\ref{sec:estimation} develops the audit, and
Section~\ref{sec:theorems} converts the sampled residuals into a
certified root error. The laws also define the learnable target class:
training favors summarizers whose realized calls are local-law
compatible, and any remaining error is structural when no such
summarizer in the chosen class can represent the target.
Section~\ref{sec:framework} states the equivalence in full.

\subsection{Learned Mergeable C-Trees}
\label{sec:v7-ctreepo-state-machine}

A C-Tree's job is to produce a root prediction that tracks the
oracle. The rest of this subsection defines the tree, oracle,
summarizer, query, equivalence relation, local context assumption,
local laws, and audit error. With these in hand, the C-TreePO problem
is fully specified, and \S\ref{sec:v7-ctreepo-correspondence}
exhibits classical mergeable summaries as worked examples.

\paragraph{The tree.}
The document is split into an ordered sequence of contiguous raw
blocks \(b_1,\dots,b_\ell\), and a binary tree \(T\) is built over
those blocks with each block sitting at one leaf. The tree is fixed
before any summarization begins; it is not itself a learned object.
Internal nodes carry an ordered (left, right) child pair, so the
composition \(s_L\oplus s_R\) at every merge is well-defined; the
unordered multiset case is recovered only when \(g\) is additionally
commutative. Balanced binary trees are the typical experimental
choice and the shape drawn in Figure~\ref{fig:v7-plain-tree}, but
the framework only requires that \(T\) be ordered binary.

\begin{definition}[C-Tree]
\label{def:v7-ctree}
A \emph{C-Tree} is a tuple \((T,\{b_v\},g,f,\rho)\) consisting of an
ordered binary tree \(T\), a raw block \(b_v\in\mathcal{B}\subseteq\mathcal{X}\) at every
leaf \(v\), a summarizer \(g\) (Definition~\ref{def:v7-summarizer}),
a query \(f\) (Definition~\ref{def:v7-query}), and a renderer
\(\rho:\mathcal{Z}\to\mathcal{X}\) used when states are not already
literal elements of \(\mathcal{X}\) and must be queried
against raw objects. Every node \(v\) represents an object
\(X_v\in\mathcal{X}\), with \(X_v=b_v\) at leaves and
\(X_v=X_L\oplus X_R\) at internal nodes. The realized state at \(v\)
is \(s_v=g(b_v)\) at leaves and \(s_v=g(s_L\oplus s_R)\) at internal
nodes; the tree prediction is
\(\widehat y_T = f(s_{\mathrm{root}})\).
\end{definition}

A run of the C-Tree has three stages. A leaf-encoding pass first
applies \(g\) to each raw block to produce the leaf states
\(s_v=g(b_v)\). A bottom-up merge pass then walks the internal nodes
and applies \(g\) to the ordered composition of the two child states;
\(\oplus\) here is whichever ordered two-child interface \(g\)
actually consumes (text concatenation, paired JSON, stacked
embeddings, or another fixed packaging chosen by \(g\)'s argument
type). The query \(f\) is finally applied once to
\(s_{\mathrm{root}}\) to produce \(\widehat y_T\), which is compared
against the oracle value \(\fstar(X_T)\). When a law compares a state
with a raw object, the answer is read through the same fixed
interface: states either already live in the queryable object space,
or are rendered through \(\rho\) before being queried. We suppress
\(\rho\) in the displays below and write \(f(z)\) directly.

The four objects in play next are the oracle the tree is trying to
match, the summarizer that produces and merges states, the query that
reads an answer off them, and the equivalence relation the local laws
will check.

\begin{definition}[Oracle]
\label{def:v7-oracle}
An \emph{oracle} is a function \(\fstar:\mathcal{X}\to\Y\) mapping
input-space objects to the task-value space.
\end{definition}

The oracle is the trusted answer for the object being queried, what an
expert rubric or gold-standard annotation workflow would produce on the
full input. It may be an expert label, a rubric value, or another
gold-standard value. Running it on every long input during training is
too expensive, which is why a tree exists.

The summarizer is the learned operator that produces leaf states from
raw blocks and merges child states into parent states.

\begin{definition}[Summarizer]
\label{def:v7-summarizer}
A \emph{summarizer} is a (possibly stochastic) map \(g\) that takes
either a leaf object \(b\in\mathcal{B}\), an already-produced state
\(z\in\mathcal{Z}\), or a fixed composition \(s_L\oplus s_R\) of two
child states, and produces a state in \(\mathcal{Z}\).
\end{definition}

The same learned operation builds a leaf summary from raw content,
builds a parent summary from two child summaries, and can re-summarize
an existing state for the C2 stability check. Every call is local:
\(g\) sees only the local input it is called on. Statements about
\(g\) are interpreted in expectation over its randomness, which
collapses to a point mass for deterministic operators.

A summarizer alone produces states; reading a task answer off any of
them is the job of the query.

\begin{definition}[Query]
\label{def:v7-query}
A \emph{query} is a (possibly stochastic) map
\(f:\mathcal{X}\to\Y\) from queryable objects to the task-value
space. Produced states are queried either because
\(\mathcal{Z}\subseteq\mathcal{X}\) by convention or because the
renderer \(\rho:\mathcal{Z}\to\mathcal{X}\) is applied first.
\end{definition}

The query turns the stored root state into the final task answer and
is applied once, at the root. Its output is what the local laws check
and what the training objective tries to match. In settings with a
closed-form oracle the query may be \(\fstar\) itself; in learned
settings it is a separate, optionally learned, calibrated component.
Once \(f\) and \(\fstar\) are in hand, a leaf law can be written
directly as \(g(b_v)\sim_f b_v\): the query gives the leaf summary and
the raw block the same task value. For internal nodes we write
\(s_v\sim_f X_v\), where \(X_v\) is the larger object represented by
the subtree. Figure~\ref{fig:v7-plain-tree} shows the resulting shape
on an eight-leaf example.

\begin{figure}[t]
\centering
\includegraphics[width=0.95\linewidth]{01_base_plain.pdf}
\caption{A C-Tree for a document partitioned into eight contiguous
blocks \(b_1,\dots,b_8\). Each block is encoded into a leaf state
\(s_i=g(b_i)\); sibling states are then merged pairwise,
\(s_{ij}=g(s_i\oplus s_j)\), up to the root. The tree shape is fixed
before any summarization begins; only the stored states change as the
run proceeds.}
\label{fig:v7-plain-tree}
\end{figure}

\paragraph{Query equivalence.}
The local laws below check \(g\)'s behavior modulo query equivalence,
not byte-equality of states. The training objective uses a task distance
\(d_Y:\Y\times\Y\to[0,\infty)\) on the same scale as the summary
error guarantee, and two objects or states are \emph{query-equivalent}
exactly when \(f\) gives them the same answer under that distance,
\[
  u\sim_f v
  \Longleftrightarrow
  d_Y\big(f(u),f(v)\big)=0,
\]
with \(\varepsilon\)-relaxation \(u\sim_{f,\varepsilon} v\) iff
\(d_Y(f(u),f(v))\le\varepsilon\). When \(f=\fstar\), query equivalence
coincides with oracle equivalence; when \(f\) uniformly recovers
\(\fstar\) within \(\varepsilon_{\mathrm{orc}}\), the two relations
agree up to a \(2\varepsilon_{\mathrm{orc}}\) slack. The audit checks
\(\sim_f\) because that is what the query can measure, and the local
laws below are stated in the \(\sim_f\) form the audit can actually
verify. Two states are query-equivalent when the query returns the
same answer on each, and the \(\varepsilon\)-relaxation allows a
stated tolerance in the task's own units.

\paragraph{Local context compatibility.}
Query equivalence as defined above is a relation on isolated pairs
\((u,v)\). The local laws below need it to compose: when one half of
an equivalent pair sits inside a larger context the tree assembles,
the equivalence must survive the surrounding material. The next
assumption is the minimal closure property that lets pair-wise
equivalences propagate up the tree, by asking the oracle to be blind
to local substitutions of equal-answer content.

\begin{assumption}[Locally context-compatible oracle]
\label{ass:context}
For every left context \(w_L\), every right context \(w_R\), and
every \(u,v\in\mathcal{X}\), \(u\sim_f v\) implies
\(w_L\oplus u\oplus w_R
  \sim_f w_L\oplus v\oplus w_R\).
\end{assumption}

Read in plain language: if two passages have the same task answer,
swapping one for the other inside the same surrounding material
leaves the answer of the whole unchanged. This is the move the tree
induction relies on, since replacing one child by an equivalent state
in the same surrounding context must then leave the parent's answer
unchanged. The assumption is local in shape (one passage is swapped)
but strong in scope (it must hold for every left and right context
the tree could place around the swap). Direct verification across all
contexts is infeasible, which is why training and the runtime audit
work as a pair: the loss pressures \(g\) toward local-law
compatibility on observed contexts during training, and the audit
gives an empirical certificate at the realized contexts a deployed
tree produces. For
content-driven tasks (distinct counting, rubric-driven coding,
document-level extraction) the assumption is plausible. It rules
out discourse-sensitive tasks where local substitution of
equivalent content changes how the surrounding text is read.

\paragraph{The three local laws.}
Three checks cover the three moves a tree run makes: encoding a leaf,
merging two child states, and re-summarizing a state that has already
been produced. If \(g\) passes all three on the realized states, the
local-to-root composition gives a certified bound on how far the
tree's root answer can drift from the oracle.

\emph{C1 (leaf sufficiency).} The leaf state preserves the answer its
raw block would receive.

\emph{C3 (merge consistency).} Merging two valid states preserves the
answer for their union.

\emph{C2 (re-summary stability).} A state can be passed through \(g\)
again without changing its answer-relevant content.

At leaves \(X_v=b_v\). At an internal node with children \(L,R\),
\(X_v=X_L\oplus X_R\). The local laws compare the state produced by
\(g\) with the object that state is supposed to represent, using query
equivalence. Their exact versions are
\begin{align}
\tag{C1}
&\forall\text{ realized leaves }v:\quad
  g(b_v)\sim_f b_v,\\
\tag{C2}
&\forall\text{ produced states }s_v:\quad g(s_v)\sim_f s_v,\\
\tag{C3}
&\forall\text{ sibling nodes }L,R\text{ with parent }v:\quad
  g(s_L\oplus s_R)\sim_f X_L\oplus X_R.
\end{align}
The leaf law uses \(b_v\) directly to avoid unnecessary notation,
since \(X_v=b_v\) at a leaf. The symbol \(X_v\) is needed for internal
nodes, where it denotes the object represented by a whole subtree. C1
is the valid-leaf condition: a leaf summary kept the answer-relevant
content of its raw block. C3 is the valid-merge condition: a parent
summary kept the answer-relevant content of the two-child span. C2 is
the reuse condition: re-applying \(g\) to a stored summary leaves its
answer stable.

Figure~\ref{fig:v7-local-laws} marks each check on the same tree from
Figure~\ref{fig:v7-plain-tree}: the C1 callout sits on a leaf encoding
(a raw block becoming a leaf state), the C3 callout sits on a sibling
merge (two child states combining into a parent), and the C2 callout
re-applies \(g\) to an already-produced state and compares the result
with the original. Each check is local:
it touches only a bounded number of nodes, and different checks can
be sampled independently. The same locality is what will let the
audit (described below) estimate the residuals from a small fraction
of nodes per run. Table~\ref{tab:v7-local-laws-sketch-link} records
the corresponding mergeable-summary obligations.

\begin{figure}[t]
\centering
\includegraphics[width=0.95\linewidth]{07_local_laws_plain.pdf}
\caption{The three local laws drawn on the same tree as
Figure~\ref{fig:v7-plain-tree}. Bolded nodes are the states the audit
samples and runs \(\fstar\) on; solid colored arrows are
\(g\)-applications; dotted arrows fan from the two compared states up
to each law's pill, which stands in for the \(\fstar\) comparator.
\textbf{C1} checks block~\(\to\)~leaf-summary (\(b_3 \to s_3\)).
\textbf{C3} checks sibling-pair~\(\to\)~parent
(\(s_5,s_6 \to s_{56}\)), comparing \(s_{56}\) against the oracle
on the raw composition \(b_5\oplus b_6\). \textbf{C2} checks re-summary by
reapplying \(g\) to \(s_{1234}\) and comparing with the ghost
\(g(s_{1234})\).
Each check is local to a bounded context and can be sampled
independently.}
\label{fig:v7-local-laws}
\end{figure}

\begin{table}[t]
\centering
\scriptsize
\setlength{\tabcolsep}{3pt}
\begin{tabular}{@{}p{0.16\linewidth}p{0.24\linewidth}p{0.25\linewidth}p{0.25\linewidth}@{}}
\toprule
Local law & C-Tree wording & Mergeable-summary wording & Answer-sufficiency reading \\
\midrule
C1: Leaf sufficiency &
The leaf state \(g(b)\) preserves the root answer that \(f\) would assign to the
raw block \(b\). &
A state produced for represented input \(X_i\) is valid for \(X_i\). &
The leaf state is answer-sufficient for its shard or leaf block: querying the
state recovers that block's task value, exactly or within the induced
\(\varepsilon\)-tolerance. \\
C3: Merge consistency &
The parent state \(g(s_L\oplus s_R)\) preserves the root answer of the object
represented by the two children. &
Merging valid states for two child objects gives a valid state for the
parent object. &
Answer-sufficiency is closed under state merge, so a tree of local merges
produces an answer-sufficient root state. \\
C2: Re-summary stability &
Re-entering an already-produced state into \(g\) does not change the root
answer. &
A valid produced summary may be reused as an input state without leaving the
same validity/query-equivalence class. &
The state remains answer-sufficient under reuse: resummarization stays in the
same query-equivalence class. \\
\bottomrule
\end{tabular}
\caption{C-Tree local laws as answer-sufficiency conditions. Read C1 as
``the leaf summary kept the answer-relevant information,'' C3 as ``the merged
summary kept the answer-relevant information for the parent,'' and C2 as
``reusing an already-produced summary is stable.''}
\label{tab:v7-local-laws-sketch-link}
\end{table}

\paragraph{Global laws and realized laws.}
The displays above are \emph{global}: they range over all leaves,
states, and sibling pairs. In deployment we only have access to the
realized inputs of one run (the actual leaf blocks
\(b_1,\dots,b_\ell\) and merge calls produced for one document). A run is
\emph{realized-valid for accuracy \(\varepsilon\)} when every realized
C1, C2, and C3 residual stays within bounds the audit can certify.
This is the form of validity the audit can actually check, since
\(\fstar\) is too expensive to run universally.

\paragraph{The certified error.}
Let \(E_{\mathrm{law}}(g,f;T)\) be the local-law error certified from
the realized C1, C2, and C3 residuals along \(T\). The tree is
\emph{certified at accuracy \(\varepsilon\)} when
\[
  E_{\mathrm{law}}(g,f;T)\le\varepsilon .
\]
When this inequality holds, the prediction
\(\widehat y_T=f(s_{\mathrm{root}})\) tracks the oracle within
\(\varepsilon\) on the realized run (with the appropriate probability
for randomized \(g\)). The audit turns many local residual checks
into one total certified error, and passing the target
\(\varepsilon\) means the stored root state is accurate enough for
this realized tree run. Section~\ref{sec:theorems} gives the formal
composition from local residuals to this root-level error.

\paragraph{Putting the pieces together.}
The full C-TreePO problem now has all its ingredients. Given an oracle
\(\fstar:\mathcal{X}\to\Y\) on the object space \(\mathcal{X}\),
assumed locally context-compatible (Assumption~\ref{ass:context}), and a
target accuracy \(\varepsilon\), C-TreePO seeks a pair \((g,f)\) with
\begin{itemize}[leftmargin=1.5em,nosep]
\item \(g\) a summarizer-merger from leaves \(b\) and child state pairs
\(s_L\oplus s_R\) to states in \(\mathcal{Z}\),
\item \(f:\mathcal X\to\Y\) a query from queryable objects to the task
value space, calibrated
against \(\fstar\),
\end{itemize}
such that on every realized tree run, the local laws C1, C2, C3 hold
under \(\sim_f\). When they do, the certified local-law error is at
most \(\varepsilon\), and the tree prediction
\(\widehat y_T = f(s_{\mathrm{root}})\) tracks \(\fstar(X_T)\) within
\(\varepsilon\). Classical mergeable summaries are the special case
where \((g,f)\) is hand-designed and the local laws hold by algebra.
Learned C-Trees establish the same obligations from training data and
a runtime audit on the realized residuals. Two questions remain.
First, how close to the local-law-compatible class can a learned
\((g,f)\) actually get in practice?
\S\ref{sec:v7-tree-comparisons} addresses this with empirical
comparisons. Second, what families of \((g,f)\) can represent the
maps the local laws ask for?
\S\ref{sec:v7-neural-operator-route} addresses this with neural
operators. Before either,
\S\ref{sec:v7-ctreepo-correspondence} exhibits classical sketches as
worked examples of the framework, and
\S\ref{sec:v7-sketch-inventory} classifies which families enter with
theorem support.

\paragraph{Estimating \(g\) and \(f\) from data.}
The classical mergeable-summary case supplies \((g,f)\) by hand: the
encode, merge, and query are engineered, the local laws hold by
algebra, and nothing has to be estimated. The C-Tree problem above
becomes interesting when \((g,f)\) must be learned for the task.
Estimating them needs access to the oracle on some manageable subset
of inputs; the form of that access is what makes the audit's certified
bound honest. A brief aside on \emph{design-based supervised learning}
sets up the form of access the audit will use.

\paragraph{An aside on design-based supervised learning.}
Design-based supervised learning fixes the sampling probabilities for
gold-standard labels in the research design. The classical estimator
under this design is the Horvitz--Thompson form, which weights each
sampled observation by the inverse of its known sampling probability
so that expectations are unbiased by construction
\citep{HorvitzThompson1952,SarndalEtAl1992,EgamiEtAl2024}.
This matters for the audit because the residuals at sampled nodes will
feed directly into a confidence-style bound on the tree's root error.
The known sampling distribution is what lets the bound be honest about
its uncertainty.

\begin{assumption}[Design-based gold supervision]
\label{ass:v7-dsl}
For each realized node \(v\) of \(T\) (leaves and internal nodes
alike), an audit policy assigns a known sampling probability
\(\pi_v\in(0,1]\). On nodes drawn under that policy, the
gold-standard oracle value \(\fstar(X_v)\) is available.
\end{assumption}

Every auditable node has a known chance of being checked, and
checked nodes receive the trusted oracle label, so the audit can
correct for which nodes were sampled. \(\pi_v\) can be a constant
rate, can vary by depth or by node type, or can adapt to intermediate
residuals; what matters is that it is set by the audit's design and
not inferred from the data. Section~\ref{sec:estimation}
discusses concrete sampling policies.

\paragraph{Observed corpora and conditional prediction risk.}
Assumption~\ref{ass:v7-dsl} only requires that sampling probabilities
be known; it does not require an online oracle, so the same setup
covers a fixed corpus the tree will be applied to, including
documents whose oracle labels we will never observe. Conditional on that finite collection and on the fitted
pipeline \((g,f)\), the realized tree nodes are the audit population.
The audit randomizes which nodes receive gold labels, logs their
inclusion probabilities \(\pi_v\), and uses inverse-probability
weighting to estimate the average oracle residual or prediction risk
for that fixed population. The unbiasedness claim is therefore
design-based: it concerns the estimated average residual over the
fixed auditable population, not the individual tree predictions
themselves. Applying the certificate to future documents, or treating
the fixed collection as representative of a broader target population,
requires a separate corpus-selection assumption. If the collection can
be treated as an i.i.d.\ or otherwise appropriate draw from the target
population, the same estimate can be read as population prediction
risk \citep{Wooldridge2010}; if the collection is nonrandomly selected,
the remaining issue is sample selection rather than node-audit design
\citep{Heckman1979,Manski1989Selection}.

\paragraph{The audit.}
With Assumption~\ref{ass:v7-dsl} in hand, the audit is the procedure
that turns sampled gold-standard labels into estimates of the realized
C1, C2, and C3 residuals. Each sampled node contributes an
inverse-probability-weighted residual; the weighted residuals are
combined into estimators of the global C1, C2, C3 quantities; those
estimators feed into the local-law error \(E_{\mathrm{law}}(g,f;T)\).
Concretely, the audit's output is a sample-based certificate that the
run is realized-valid at accuracy \(\varepsilon\), with whatever
confidence the sampling design supports.
Section~\ref{sec:estimation} develops the audit in full;
Section~\ref{sec:theorems} compiles the residuals into the formal
bound. Classical mergeable summaries have zero residuals by
construction; learned C-Trees require the audit.

\subsection{Correspondence with Classical Sketches}
\label{sec:v7-ctreepo-correspondence}

The clean comparison is same-tree first. Fix a leaf partition and a
binary tree \(T\). A full classical mergeable summary can be run on
that same \(T\): encode at leaves, merge states at internal nodes, and
query the root. Its original Agarwal-style validity and mergeability
obligations give the C1/C3 part of the C-TreePO local laws. If the
scheme also supplies stable re-entry for produced states, it supplies
C2 as well. Conversely, exact C-TreePO local laws imply that the root
summary distribution is supported on oracle-valid summaries for the
represented input. Approximate local laws give the corresponding
\(\varepsilon\)- or expected-validity statement, not exact state
equality. What classical sketches earn by algebra, learned C-Trees
must earn by audit.

\paragraph{Forward direction.}
A classical mergeable summary supplies a same-topology C-Tree by using
its \(\mathrm{encode}\) at leaves, its \(\oplus\) at internal nodes,
and its \(\mathrm{query}\) at the root. C1 and C3 then hold by the
sketch's own validity and merge-closure obligations. C2 is the extra
law C-TreePO requires beyond the baseline Agarwal obligations: the
stable-reentry condition needed when produced states are fed back into
the summarizer.

\paragraph{Reverse direction.}
If C1 and C3 hold exactly on \(T\), the one-pass root distribution is
supported on states \(z\) with \(d_Y(\fstar(z),\fstar(S(T)))=0\). With
C2 added, the same supportwise validity holds for multi-round \(Z_R\)
reductions, since C2 ensures that re-entering a produced summary
through \(g\) preserves its oracle value. Recovering a fully classical
state-level summary additionally requires a validity relation, a
query, and (for Agarwal's sized notation) an external size profile
\(k(|D|,\varepsilon)\); without these the correspondence is
behavioral at the level of task answers, with byte-level state equality
only when the chosen serialization or codec is invertible. The
following proposition states the resulting equivalence formally.

\begin{proposition}[Exact query equivalence]
\label{prop:exact-query-equivalence}
Fix a C-Tree \(T\) with \(S(T)=x\) and a classical state-merge tree
\(t\) whose stream data represents the same input, so that
\(\mathrm{oracle}(\mathrm{data}(t))=\fstar(x)\). Suppose C1, C2, and
C3 hold exactly for \(g,T,\fstar\), and suppose the classical summary
satisfies build validity, merge validity, and valid-state query
correctness. Then for every realized C-Tree output
\(z\in\mathrm{supp}(Z_R(g,x,T))\), \(R\ge 1\),
\[
  \fstar(z)
  =
  \mathrm{query}\!\left(\mathrm{eval}_t(\mathrm{build},\mathrm{merge})\right)
  =
  \mathrm{oracle}(\mathrm{data}(t)).
\]
Write
\[
  \mathcal{S}_{\mathrm{Ag}}(D,\varepsilon)
  =
  \{s:\mathrm{Valid}_{\varepsilon}(D,s)\}.
\]
If \(D\) represents the same input \(x\), then every realized C-Tree
output \(z\) and every
\(s\in\mathcal{S}_{\mathrm{Ag}}(D,\varepsilon)\) satisfy
\[
  \fstar(z)=\mathrm{query}(s)=\mathrm{oracle}(D).
\]
If the C-Tree state space is identified with the Agarwal state space by
an explicit validity bridge, then
\(\mathrm{supp}(Z_R(g,x,T))\subseteq
\mathcal{S}_{\mathrm{Ag}}(D,\varepsilon)\).
\end{proposition}

This is the state-level nesting claim in plain form: a C-Tree with
exact local laws and a matching classical summary gives the same root
answer as the classical merge tree. The equivalence is at the
level of query values; literal state equality requires an additional
invertible codec, and sized containment requires the external
\(k(|D|,\varepsilon)\) profile.

\begin{proposition}[Reduction to a classical sketch]
\label{prop:mergeable-reduction}
Under Assumption~\ref{ass:context}, suppose \(g\) is deterministic
and satisfies the global preservation laws
\[
\text{(A1)}\ \forall s\in\mathcal{Z},\ g(s)\sim_f s,\qquad
\text{(A2)}\ \forall u,v\in\mathcal{X},\
  g(g(u)\oplus g(v))\sim_f u\oplus v.
\]
Let \((\mathrm{encode},\oplus_B,\mathrm{query})\) be a bounded
mergeable-summary scheme \citep{Agarwal2013MergeableTODS}. If \(g\)
serializes the sketch state, with \(g(x)=\mathrm{encode}(x)\) on
raw shards/leaf blocks and \(g(z_1\oplus z_2)=z_1\oplus_B z_2\) on
serialized states, and \(f(z,q)=\mathrm{query}(z,q)\), then the C-Tree
and the classical summary produce the same state-level merge and the
same root answer.
\end{proposition}

This proposition gives the forward construction: a bounded mergeable
summary supplies a C-Tree instantiation by reusing its components.
The classical encode and merge become the summarizer \(g\); the
classical query becomes the C-Tree query \(f\). Bounded sketch states
compose by \(\oplus_B\), \(f\) is applied only after the state has
been merged, and A1, A2 give the C-Tree-side preservation. Full proof
in Appendix~\ref{app:proof-mergeable}.

\paragraph{Randomized sketches.}
The same statement holds with the original sketch's success
probability \(p\): if the classical root state is valid with
probability at least \(p\), the C-Tree instantiation satisfies the
query guarantee with probability at least \(p\) as well.

\paragraph{Special case: strict oracle homomorphism.}
Under the same A1 and A2, suppose additionally that the task value
itself admits a merge: there exists
\(m_{\Y}:\Y\times\Y\to\Y\) with
\(m_{\Y}(\fstar(u),\fstar(v))=\fstar(u\oplus v)\) for all \(u,v\).
Then \(\fstar\) is a homomorphic query of the induced state.
Merging summaries and reading out gives the same oracle value as
reading the raw concatenation, and \(\oplus_g\) is associative
modulo \(\sim_f\). Most C-Tree tasks have no such scalar merge; the
state-level reduction above is the operational route.

\paragraph{Schedule invariance.}
The same leaves admit many merge trees, and the local laws do not
care which one is used once they hold on the realized calls. Whenever
C1 and C3 hold exactly on two ordered trees with the same leaves, both
root distributions have zero oracle distortion, so the two schedules
are equivalent at the query level. With approximate laws, the
corresponding statement is an expected or certified error bound for
each certified schedule; it is not a claim of identical internal
states.

\paragraph{Scalar-query caveat.}
The tree merges states, not scalar answers. Distinct count is the
running example: knowing the distinct count in leaf block \(A\) and in
leaf block \(B\) does not determine the count in \(A\uplus B\). A
state-level summary that retains hash fingerprints can merge; two
scalar counts cannot.

\subsection{Sketch Families and Claim Status}
\label{sec:v7-sketch-inventory}

The correspondence above tells us \emph{when} a classical sketch
instantiates a C-Tree directly; the next move is classifying
\emph{which} concrete sketch families enter with theorem support. Different families enter with different
backing: some carry a formal or Agarwal-style obligation we can rest the
local laws on, others enter only as empirical references, and a few
appear as adjacent geometric background that should not be attributed
to the same source. Table~\ref{tab:v7-sketch-inventory} records the
classification. The distinction matters for citation honesty:
\citet{Agarwal2013MergeableTODS} lacks an \(\varepsilon\)-kernel
theorem, so \(\varepsilon\)-kernel geometry is adjacent background
unless cited to a separate source.

\begin{table}[t]
\centering
\scriptsize
\setlength{\tabcolsep}{3pt}
\begin{tabular}{@{}p{0.16\linewidth}p{0.26\linewidth}p{0.18\linewidth}p{0.27\linewidth}@{}}
\toprule
Family & Description & Theorem status & What the C-Tree uses \\
\midrule
HLL & Estimates the number of distinct items without storing the items
themselves; the state is a vector of registers. & theorem-backed + empirical & Register-state summary, pointwise-max merge, and final estimator query value; the tree and flat reductions agree when they use the same state representation. \\
Count-Min & Estimates item frequencies with hashed counter tables; the query
asks for the approximate count of one item. & theorem-backed + empirical & Additive table summary and point-frequency query value; mergeable but not idempotent. \\
Misra--Gries / SpaceSaving & Find heavy hitters: items whose frequencies are
large enough that they should survive compression. & theorem-backed interfaces & Capacity, mass, and error-envelope arguments for Misra--Gries; SpaceSaving follows by the standard min-counter correspondence. \\
GK quantiles & Maintains a compressed ordered summary so approximate ranks and
quantiles can be answered from a stream. & theorem-backed one-way & Sequential/incremental quantile interface; not claimed as arbitrary full state-state mergeability. \\
KLL quantiles & Randomized quantile sketch that compacts weighted samples and is
designed for mergeable quantile states. & randomized theorem-backed interface & Randomized sized mergeable query-sketch interface with a weighted-mass invariant. \\
Interval/range constructions & Keep small subsets or colored states that approximate many interval or range counts at once. & theorem-backed scaffolds + conditional size-reduction constructions & Same-weight interval concentration, finite VC/Sauer--Shelah trace layer, and low-discrepancy coloring route when supplied by a certificate. \\
Markov and bigram sketches & Store the count information needed to answer simple sequence-model queries, such as transition or bigram statistics. & theorem-backed local-law examples & Task-specific answer-sufficient states showing the same summarize/merge/query pattern outside classical data-stream queries. \\
CPC, Theta/KMV, Frequent Items, classic quantiles, REQ, t-digest, Tuple, VarOpt & Additional library sketch families used to compare tree reductions against familiar implementations. & empirical-only here & Empirical comparisons only; no theorem is claimed for these additional variants. \\
\(\varepsilon\)-kernel geometry & Geometric approximation family for preserving directional or shape information, separate from the mergeable-summary results of \citet{Agarwal2013MergeableTODS}. & adjacent/background here & Do not attribute to that paper without a separate source. \\
\bottomrule
\end{tabular}
\caption{Mergeable-sketch claim status. Rows record the claims used
here, not the full external literature on each sketch family.}
\label{tab:v7-sketch-inventory}
\end{table}

\subsection{Tree Comparisons Against Classical and Learned Sketches}
\label{sec:v7-tree-comparisons}

The empirical comparisons use the correspondence in three ways. The
classical anchor here is HyperLogLog
\citep{FlajoletEtAl2007,HeuleEtAl2013}, the standard register-vector
sketch for distinct counts.

\paragraph{Classical state equality.}
Classical HLL register vectors agree with the flat-sketch register
vector on every tested tree shape; this is the tightest possible parity
check, since the state is explicit. With a C-Tree instantiated using
HLL's register state and pointwise-max merge, the tree-reduced register
vector equals the flat register vector for every tree shape we tested.
Alternative HLL variants can agree at the estimate level even when
their internal representation changes along different construction
paths.

\begin{figure}[t]
\centering
\includegraphics[width=0.92\linewidth]{classical_sketches_hll_leaf_size.pdf}
\caption{\textbf{Official HLL over leaf size.}
The exact HLL anchor uses the native register state and the fixed HLL
estimator \(f^\star\). Because the state merge is register-wise max, the
tree reduction agrees with the corresponding flat sketch at the state level;
the curve therefore measures the ordinary HLL estimation floor. Learned-merge
error is outside this anchor. Learned HLL-register rows, when included, are
interpreted as recovery of this supplied state/query pair.}
\label{fig:v7-hll-parity-curves}
\end{figure}

\paragraph{Learned state with fixed versus learned query.}
The learned HLL rows ask whether a neural \(g\) can recover a classical
state-and-estimator target. With a fixed classical HLL estimator as
\(f\), merge errors must land in the exact register geometry and can
compound with tree depth. With learned \(g+f\), the system can choose a
different latent state that preserves the same answer. The practical
lesson for semantic C-Tree tasks: exact recovery of a reference
sketch imposes a real query-shape constraint; task preservation alone
may be better served by a jointly learned state and query.

\paragraph{Broad sketch suite.}
The same tree-reduction protocol applies to classical sketches and
learned companions across distinct-counting, frequency, quantile, set,
and sampling-style queries.
Table~\ref{tab:v7-classical-sketches-compact} reports representative
rows; larger grids are omitted to keep the discussion focused.

For sketches lacking an exact locally mergeable law in the state class
we train, the learned rows are approximations. The optimizer searches the
chosen mergeable-summary class for a state/query pair whose tree
reduction best matches the empirical oracle under the training
distribution, with no claim to recover the library's internal update
rule. Exact recovery occurs only when the target law has a
representative in that class; otherwise the residual gap is a
structural error beyond optimization failure.

\begin{figure}[t]
\centering
\includegraphics[width=0.92\linewidth]{classical_sketches_summary.pdf}
\caption{\textbf{Broad mergeable-sketch approximation view.}
The x-axis is leaf size in tokens, with the top axis showing the induced
number of leaves per document. Official/oracle rows are the reference
tree-reduction implementations. Learned non-HLL rows, when present, report
the excess structural error relative to the reference oracle, with no claim
that the exact library law has been recovered.}
\label{fig:v7-classical-sketches-summary}
\end{figure}

\begin{table}[t]
\centering
\scriptsize
\setlength{\tabcolsep}{4pt}
\begin{tabular}{@{}p{0.22\linewidth}p{0.32\linewidth}p{0.34\linewidth}@{}}
\toprule
Task family & Classical comparison & Learned-tree comparison \\
\midrule
Distinct counting & HLL gives the canonical state-level parity check:
state merge matches flat reduction when the same register representation is
used. & Learned HLL-like states reveal the cost of requiring exact reference
state recovery versus allowing a learned query. \\
Frequency estimation & Count-Min has additive table states and point-query
values. & Learned companions test whether the tree can preserve frequency
information under repeated merges. \\
Quantiles and ranks & GK and KLL represent the deterministic one-way and
randomized mergeable quantile regimes. & Learned companions test whether a
answer-sufficient semantic state can preserve rank information well enough for
the root query. \\
Set and sampling-style summaries & Additional classical sketches serve as
empirical baselines across set, sampling, and weighted-summary tasks. &
Learned rows are treated as empirical evidence only unless they satisfy the
state-level local laws above. \\
\bottomrule
\end{tabular}
\caption{\textbf{Broad-sketch comparison.} Classical rows are empirical
references; theorem status is the separate classification in
Table~\ref{tab:v7-sketch-inventory}.}
\label{tab:v7-classical-sketches-compact}
\end{table}

\paragraph{Learning an answer-sufficient state from oracle feedback.}
When the state dimension \(m\) can carry the \(r\)-dimensional
answer-sufficient count vector, the learned sketch approaches the
hand-designed state; when \(m<r\), the error floor is structural. The
controlled diagnostic asks the oracle for an \(r\)-vector of per-type
counts, so a state with dimension \(m<r\) is structurally insufficient
and the learned tree cannot close the gap by training alone. The
learned tree converges when \(m\ge r\) and plateaus when \(m<r\),
matching the hand-designed sketch boundary. The Manifesto setting uses
the same logic with a semantic answer-sufficient state.

\begin{figure}[t]
\centering
\includegraphics[width=0.88\linewidth]{learned_sketch_leaf_size_diagnostic.pdf}
\caption{\textbf{Learning a local-law-compatible approximation from oracle feedback.} When
the state dimension \(m\) can carry the answer-sufficient state, the learned
tree approaches the hand-designed boundary. When the target behavior is
outside the chosen mergeable class, the learned row estimates the closest
local-law-compatible approximation under the empirical loss, and the
remaining gap is structural.}
\label{fig:v7-learned-sketch-curves}
\end{figure}

\paragraph{Non-commutative oracles.}
Sketches with associative, commutative, idempotent merges are friendly
first examples. Most interesting C-Tree oracles are not
register-like. The Markov-changepoint example of
Section~\ref{sec:markov-walkthrough} counts boundary transitions
between adjacent colored regions, a cross-span interaction no
commutative register-wise merge can recover from leaves alone. The
merge has to carry boundary information forward, which is what
motivates the Fourier neural operator described in
Section~\ref{sec:fno-primer}. The audit-and-local-law picture is the
same; the architecture of \(g\) has to be richer.

\subsection{Neural Operators as a Representation Route}
\label{sec:v7-neural-operator-route}

The local laws say what \(g\) must satisfy on realized states. The
architecture of \(g\) determines which summarizers are reachable in the
first place. A classical mergeable summary's \(g\) is hand-written: a
fixed-arity encode and a fixed-arity merge, both engineered from the
algebra of one query. A neural \(g\) is a parameterized family trained
under local-law pressure and then audited on the realized calls. The
question is what families to use.

A \emph{neural operator} is a network designed to act on structured
inputs (functions, fields, or high-dimensional states); a flat token
sequence is only one possible representation
\citep{Kovachki23NeuralOperator}. The Fourier neural operator used in
the Markov example treats a leaf's hidden state as a function on a
discrete spatial axis, applies layered Fourier-domain transforms, and
produces another such state. Transformer-as-operator blocks fall in the
same class. The key property is that the family is rich enough to
represent ideal leaf-summary and state-merge maps on the finite
collection of realized calls a tree generates, while small enough to
train. If the state space is represented in the operator domain, an
appropriate neural-operator class can approximate the ideal \(g\) on
that realized call set.

\begin{proposition}[Fixed-\(\varepsilon\) learned recovery]
\label{prop:fixed-epsilon-no-recovery}
Fix a target accuracy \(\varepsilon\), a tree \(T\), and \(R\ge1\).
Suppose an ideal summarizer \(g^\star\) in a neural-operator class
\(\mathcal C\) satisfies the exact local laws on \(T\). Suppose a
realized neural operator \(g_\theta\in\mathcal C\) approximates
\(g^\star\) on the realized leaf, merge, and re-entry call sets, and
that the approximation transfers to local-law budgets
\(\varepsilon_{\mathrm{leaf}},\varepsilon_{\mathrm{merge}},
\varepsilon_{\mathrm{idemp}}\). If
\[
  \varepsilon_{\mathrm{leaf}}
  +\varepsilon_{\mathrm{merge}}
  +(R-1)\varepsilon_{\mathrm{idemp}}
  \le \varepsilon ,
\]
then the root distribution of the learned tree is expected-\(\varepsilon\)
valid:
\[
  \mathbb E_{z\sim Z_R(g_\theta,x,T)}
  d_Y(\fstar(z),\fstar(x))
  \le \varepsilon .
\]
The same statement applies to the FNO route once the finite-mode/FNO
approximation certificate supplies the same realized-call bounds. If the
laws are checked through a calibrated learned query \(\widehat f\), the
condition becomes
\[
  \varepsilon_{\mathrm{leaf}}
  +\varepsilon_{\mathrm{merge}}
  +(R-1)\varepsilon_{\mathrm{idemp}}
  +2\varepsilon_{\mathrm{orc}}
  \le \varepsilon ,
\]
where \(\varepsilon_{\mathrm{orc}}\) is the uniform oracle-recovery error of
\(\widehat f\) against \(\fstar\).
\end{proposition}

This proposition is the fixed-target certification route:
approximation theory supplies a representable \(g_\theta\) on the
realized calls, optimization and audit supply the local-law budget, and
the theorem converts that budget into expected-\(\varepsilon\) validity.
In existential form, any NO or FNO learning result that returns such a
certified \(g_\theta\) gives an \(\varepsilon\)-good learned C-Tree
summary for the fixed \(T,x,R\).

A useful way to state the learned-oracle assumption is separate from the
summarizer assumption. If the NO/FNO query \(\widehat f\) recovers the
true oracle uniformly,
\[
  \sup_x d_Y(\widehat f(x),\fstar(x)) \le \varepsilon_{\mathrm{orc}},
\]
and the learned tree has local-law budget \(E_{\mathrm{law},R}\) measured
through \(\widehat f\), then the same tree has true-oracle expected
distortion
\[
  \mathbb E_{z\sim Z_R(g_\theta,x,T)}
  d_Y(\fstar(z),\fstar(x))
  \le E_{\mathrm{law},R}+2\varepsilon_{\mathrm{orc}}.
\]
Thus a fixed target \(\varepsilon\) is obtained whenever
\(E_{\mathrm{law},R}+2\varepsilon_{\mathrm{orc}}\le\varepsilon\).

Approximation and certification are separate steps. A flexible
architecture may learn a useful summarizer; the local laws and audit
tell us whether the learned state behaves like a mergeable summary for
the task. In the differentiable neural-operator setting, the empirical
objective has the form
\[
  \mathcal{L}
  =
  \alpha_{\mathrm{cal}}\mathrm{calibration\_loss}(\widehat f)
  +\alpha_{\mathrm{gold}}\mathrm{gold\_loss}(\widehat f,g)
  +\lambda\,E_{\mathrm{law},R}(\widehat f,g).
\]
The gold loss measures the distance from the tree's final prediction
to the available gold label. The local-law term is the weighted
empirical C1/C3/C2 residual. The summarizer \(g\) enters through the
tree states it produces. The learned query \(\widehat f\) enters through
calibration, gold-label prediction, and query-level local residuals.
The coefficient \(\lambda\) is a training hyperparameter; it is not an
error unit. Certification is the separate statement that the total
certified error is at most the target \(\varepsilon\).

When local oracle labels are not available, as in the LLM setting, we
use a teacher-first version of the same objective. The first pass uses
raw documents or raw concatenations to train a learned query
\(\widehat f\) from global labels. The next pass fixes or selects this
teacher and uses it to define local residuals: C1 compares
\(\widehat f(g(b))\) with \(\widehat f(b)\), C3 compares
\(\widehat f(g(z_L\oplus z_R))\) with
\(\widehat f(X_L\oplus X_R)\), and C2 compares
  \(\widehat f(g(z))\) with \(\widehat f(z)\). The final claim then
  adds the teacher oracle-recovery slack to the audited local-law error.

The slack comes from a simple residual transfer. For any local
comparison \(a,b\),
\[
  d_Y(f^\star(a),f^\star(b))
  \le
  d_Y(\widehat f(a),\widehat f(b))
  +d_Y(\widehat f(a),f^\star(a))
  +d_Y(\widehat f(b),f^\star(b)).
\]
If the teacher has uniform oracle-recovery error at most
\(\varepsilon_{\mathrm{orc}}\), each
local residual measured through \(\widehat f\) transfers to the
corresponding true-oracle residual with at most
\(2\varepsilon_{\mathrm{orc}}\)
additive slack. Thus the componentwise true local-law error is bounded by
\[
  E_{\mathrm{law}}(g,f^\star;T)
  \le
  E_{\mathrm{law}}(g,\widehat f;T)+2(R+1)\varepsilon_{\mathrm{orc}}
\]
for the unweighted C1/C3/C2 composition with \(R-1\) re-summary
terms. For the final root prediction we use the sharper end-to-end
bound below, which transfers the root comparison once and therefore
pays only \(2\varepsilon_{\mathrm{orc}}\).

A compact way to state the learned guarantee is to expand the DSL/IPW
certificate rather than hide it in a new symbol. Let
\(\widehat E_{\mathrm{law},R}^{\mathrm{IPW}}\) be the HT/H\'ajek estimate
of the C1/C3/C2 residual composition for the learned query
\(\widehat f\), let \(\widehat{\mathrm{SE}}_{\mathrm{IPW}}\) be its
reported DSL/IPW standard error, and let \(B_{\mathrm{cal}}\) be the
held-out query/oracle calibration margin. If \(\widehat f\) uniformly
recovers \(\fstar\) within \(\varepsilon_{\mathrm{orc}}\), the practical
certificate is
\[
\begin{aligned}
  \mathbb E_{z\sim Z_R(g_\theta,x,T)}
  d_Y(\fstar(z),\fstar(x))
  &\le
  \underbrace{\widehat E_{\mathrm{law},R}^{\mathrm{IPW}}}
    _{\text{sampled local-law estimate}}
  +
  \underbrace{z\,\widehat{\mathrm{SE}}_{\mathrm{IPW}}}
    _{\text{DSL/IPW uncertainty}} \\
  &\quad+
  \underbrace{B_{\mathrm{cal}}}
    _{\text{calibration margin}}
  +
  \underbrace{2\varepsilon_{\mathrm{orc}}}
    _{\text{oracle recovery}} .
\end{aligned}
\]
Thus the learned tree is certified for \(\fstar\) at tolerance
\(\varepsilon\) whenever the right-hand side is at most \(\varepsilon\).
The factor \(2\varepsilon_{\mathrm{orc}}\) is the two-sided cost of comparing
through the calibrated query \(\widehat f\). Under the DSL/IPW sampling
assumptions, the IPW estimate is unbiased or consistent for the realized
local-law estimand; increasing effective sample size shrinks the uncertainty
term, not automatically the learned residual itself. A neural-operator
architecture is covered when the realized state machine lies in, or is certified
near, the local-law-compatible class.

The learned query has its own oracle-recovery error, and the
certificate for the learned tree must cover both the audited tree
residuals and that recovery slack. This gives three readings of
the same error formula. With only root labels, the local laws are
measured through a calibrated teacher \(\widehat f\), so the
certificate pays the teacher slack. With online oracle calls and a
genuinely mergeable oracle, the local laws can be checked directly
against \(f^\star\), and the exact mergeable witness has zero law
error. With online oracle calls but a non-mergeable target, the
procedure is no longer recovering \(f^\star\) exactly; it is fitting
\(g_\theta\) to the best available tradeoff between root accuracy and
local-law error in the learned class. The remaining root error
\(E_{\mathrm{root}}(g_\theta,f^\star;T)\), meaning the true-oracle
error at the tree root, and
\(E_{\mathrm{law}}(g_\theta,f^\star;R)\) are structural when no exact
mergeable representative exists.

\subsection{Consequences}
\label{sec:v7-what-established}

The preceding discussion establishes the following chain:
\begin{enumerate}[leftmargin=1.5em]
\item Classical mergeable summaries summarize leaves into states, merge
states, and query the root.
\item C-Trees nest these state-level objects through relational
state-level trees; scalar child-answer merging is a separate stronger
assumption.
\item Randomized mergeable summaries nest with the same probability as
their root-validity event.
\item Exact local laws induce a mergeable state scheme on the realized
range of \(g\); known sketches supply direct witnesses for those laws.
\item Neural operators supply candidate summarizers for leaf and merge
calls. The C-Tree local laws plus a runtime audit replace the classical
sketch's once-and-for-all algebraic proof with sampled residuals on a
deployed run.
\item The empirical suite checks both sides: exact collapse to
classical HLL when the state is fixed, and learned recovery or
structural failure when the state is trained from oracle feedback.
\end{enumerate}

The local-law-compatible class described by C-Trees captures the
classical mergeable-summary setting and extends it to learned semantic
states such as the Manifesto application in
Section~\ref{sec:min-manifesto-results}.
